% Generado por src/comparativa.py — no editar a mano, se regenera.
\begin{table}[htbp]
    \centering
    {\footnotesize
    \begin{tabular}{@{}lrrrrrr@{}}
        \toprule
        \textbf{Observable} & \textbf{n (cobertura)} & \textbf{Ridge} & \textbf{RF} & \textbf{SAGE} & \textbf{GIN} & \textbf{GATv2} \\
        \midrule
        media de $Z$ & 1986 (50\%) & 1,0000 & 1,0000 & 1,0000 & 1,0000 & 1,0000 \\
        media de $X$ & 3393 (85\%) & 1,0000 & 1,0000 & 1,0000 & 1,0000 & 1,0000 \\
        media de $Y$ & 1096 (27\%) & 1,0000 & 1,0000 & 1,0000 & 1,0000 & 1,0000 \\
        paridad & 1825 (46\%) & 1,0000 & 1,0000 & 1,0000 & 1,0000 & 1,0000 \\
        corr. de vecinos & 3369 (84\%) & 1,0000 & 1,0000 & 1,0000 & 1,0000 & 1,0000 \\
        \bottomrule
    \end{tabular}}

    \caption{Acierto de signo del observable mitigado. Conjunto de test, por observable. $n$ es el número de casillas evaluables. Ridge = Ridge de 34 variables, RF = Random Forest; SAGE, GIN y GATv2 son las tres convoluciones del MPNN.  Son métricas post-ejecución: se reportan, pero la comparativa la decide el $R^2$ sobre $f$.}
    \label{tab:mitigacion_signo}

    {\footnotesize Fuente: elaboración propia.}
\end{table}
